\subsection{Overview of the STAR protocol}
\label{subsec-STAR}
Our proposed construction is inspired by the STAR protocol, which we briefly recall here.
In STAR, a client with value $v$ computes a $\numheavyhit$-out-of-$\numclients$ secret sharing of their value $\share{v} = (\sharei{v}{1}, \ldots, \sharei{v}{\numclients})$, samples an index $i$, and sends $\sharei{v}{i}$ to the Aggregation Server.
When $\ell$ clients report the same value $\val$, the Aggregation Server learns $\sharei{\val}{i_1}, \ldots, \sharei{\val}{i_\ell}$ for $\ell$ distinct indices $\set{i_1, \ldots, i_\ell}$.
If $\ell < \numheavyhit$, then $\val$ is not a heavy-hitter and the privacy of the $\numheavyhit$-out-of-$\numclients$ secret sharing ensures that it remains private.
On the other hand, when $\ell > \numheavyhit$, the Aggregation Server possesses more than $\numheavyhit$ shares with which it can reconstruct the heavy-hitter $\val$.

Unfortunately, the protocol as described does not ensure correctness, as the standard reconstruction algorithm for a $\numheavyhit$-out-of-$\numclients$ threshold secret sharing requires at least $\numheavyhit$ shares from the \emph{same} sharing instance.
Specifically, this leads to two issues with the above protocol:
\begin{enumerate}
	\item Clients with the same input do not report shares from the same sharing instance.
	\item Even if they did, the Aggregation Server receives only a list of shares and cannot determine which subset corresponds to the same value, making reconstruction infeasible.
\end{enumerate}
STAR addresses the first issue by having the Randomness Server coordinate the randomness, so that clients with the same input generate correlated shares from the same sharing instance.
In more detail, each client with input $v$ obliviously obtains randomness $r_v$ associated with $v$---so that the Randomness Server does not learn $v$, and the client does not learn the randomness for other inputs\footnote{This can be done using an Oblivious Pseudorandom Function (OPRF) protocol.}---and then uses $r_v$ to deterministically derive the randomness used to secret share $v$.
Each client then sends a share at a distinct index to the Aggregation Server.

To resolve the second issue, STAR requires each client to append a tag $t_v$, also derived from $r_v$, to their report.
Each client sends a report of the form $\report_i = (\sharei{v}{i}, t_v)$ to the Aggregation Server, which then combines $\numheavyhit$ reports with the same tag to reconstruct the heavy-hitters.
The inclusion of the tags $t_v$ in the STAR protocol leads to its core problem: leaking the \emph{entire} histogram of client values.

\subsection{Our Approach}

\subsubsection{Reducing Leakage}

The first contribution of this work is an approach to reducing the STAR leakage.
In contrast to POPSTAR, which still allows rare values to have their frequencies leaked, we endeavor to provide complete privacy for clients with low frequency reports.

To do this, we remove the report tags entirely, noting that the problem of reconstructing secrets without tags can be solved using Multi-Dealer Secret Sharing (MDSS), originally described in \autocite{\citeairtag}.
MDSS extends traditional secret sharing by allowing for the simultaneous recovery of multiple secrets from a pool of \emph{unlabeled} shares, that may be of unrelated secrets.
In our setting this ensures that heavy-hitters can be extracted from a set of reports without the need to explicitly label those for a specific input value, and without revealing any information about non-heavy-hitter values. In principle, MDSS offers us a drop-in solution to the tag leakage problem:
clients query the randomness server to obtain the secret-sharing parameters corresponding to their input, evaluate a single share, then send that share to the Aggregation Server without a tag.
The Aggregation Server then runs the MDSS reconstruction algorithm over the entire pool of shares to learn the heavy hitters.


While MDSS can theoretically work as a drop-in solution, making this approach work efficiently in practice is non-trivial.
One important concern is that existing MDSS constructions add a gap (or \emph{ramp}) between the number of shares required for secret recovery and the number of shares that are guaranteed to protect privacy (throughout this work we refer to these thresholds as $\threshold$ and $\pd$ respectively). This ramp can be minimized by increasing the size of the MDSS shares and the reconstruction runtime, although it cannot be practically eliminated.
This requires some careful balancing of parameters, as well as a number of optimizations to ensure a practical level of efficiency. We describe these in \hme{ADD BACK}.
% \cref{sec:star-implementation}.



\subsubsection{Preventing Spying}
We next turn towards preventing the spying attack, where the Aggregation Server can identify all shares of a value of choice after a single query to the Randomness Server.

Removing the tags associated with the reports takes a step towards alleviating this concern, but does not address the problem on its own.
In STAR, clients do not just learn a share of their value, they learn the entire \emph{sharing}: $(\sharei{\val}{1}, \dots, \sharei{\val}{\numclients})$.
Given this, the Aggregation Server can still perform the spying attack by first submitting the target value $v$ to the Randomness Server to learn $(\sharei{\val}{1}, \dots, \sharei{\val}{\numclients})$, and then checking if any report is one of the $\sharei{v}{i}$.

Addressing this problem requires changing how clients interact with the Randomness Server. At a high level, we require a 2-party computation (2PC) protocol that deterministically maps a client input $v$ to a seed value $r_v$, uses it to compute a secret sharing, and returns a single share to the client, all while keeping the sharing hidden from both parties. While we could achieve this using generic 2PC techniques, we would like to keep the Randomness Server interaction as lightweight as possible and generic 2PC introduces a significant overhead.

To overcome this hurdle we exploit the structure of MDSS shares.
An MDSS share is a point on a tuple of polynomials: $(x, p_1(x), \dots, p_{\np}(x))$.
Critically, the share algorithm consists only of sampling these polynomials uniformly at random, and then evaluating them on the same random point.
We therefore give a protocol for {\em doubly oblivious polynomial evaluation} (DOPE) that maps an input to a seed, derives a polynomial, and evaluates that polynomial on a given point, all while keeping both parties ignorant of the intermediate steps.


To construct such a protocol, we use an additively homomorphic public key encryption scheme that allows for the oblivious sampling of ciphertexts.
In such a scheme, parties can use a public key to sample valid ciphertexts without knowing the underlying plaintexts.

Given this, our protocol begins with a client hashing their value $v$ to obtain $r_v$. They then use the public key of the Randomness Server and $r_v$ to obliviously sample $\pd$ ciphertexts $(c_1, \dots, c_{\pd})$, with corresponding plaintexts $(m_1,\dots,m_{\pd})$.
The client then computes the polynomial evaluation $y = c_1x + c_2x^2 + \dots c_{\pd}x^{\pd}$ for a value $x$.
By the homomorphic property of the encryption scheme the result is the encryption of $m_1x + m_2x^2 + \dots + m_{\pd}x^{\pd}$.
Finally, the client samples a mask $r^*$, computes its encryption $c^*$, and sends $y + c^*$ to the Randomness Server.
The Randomness Server decrypts the ciphertext and sends the result to the client, who removes the mask.
At the end of the protocol, the client is left with only a single point on an unknown polynomial, and the Randomness Server learns nothing due to the single value it decrypts being masked.
Knowledge of the value $r_v$ is no longer useful to a corrupted client, as it requires interaction with the Randomness Server to learn any information about additional shares.
Critically, to learn the polynomial associated with $v$ (and thereby be able to identify other reports for $v$), a colluding Aggregation Server would need to issue $\pd + 1$ queries to the Randomness Server.

There are several encryption schemes that are both $(1)$ additively homomorphic and $(2)$ allow for oblivious ciphertext sampling: these include the Paillier \cite{EC:Paillier99} and Goldwasser-Micali \cite{STOC:GolMic82} schemes. We ultimately elect to use the Benaloh encryption scheme \cite{THESIS:Ben87,SAC:Ben94}, as its plaintext space aligns well with the share space required by MDSS. For more information see Section \hme{Add back}.
% \ref{sec:dope}.

For our scheme, the interaction between the Randomness Server and client in the DOPE protocol remains lightweight. Notably, for each polynomial, the client's DOPE message is a single Benaloh ciphertext, and the Server only performs a single decryption. Concretely, for the number of polynomials used in our work, both the client and Randomness Server computation can be performed in approximately $1$ second. Total communication is in the range of $50$ kilobytes.
More detailed benchmarks can be found in \hme{ADD BACK}.
% \cref{sec:star-experiments}.


\subsection{Comparison to POPSTAR}
\label{subsec:popstar-leakage}
\label{subsec-comparison-to-prior-work}

POPSTAR \autocite{\citepopstar} is the state-of-the art for computing private heavy hitters in the STAR model, and we briefly describe the leakage induced by the scheme here.
While the original STAR protocol reveals the client histogram, allowing the Aggregation Server to know exactly which reports are of the same value and which are not, the leakage in POPSTAR is akin to a {\em partial} revealing of the histogram.
The Aggregation Server may learn the exact counts for some of the non-heavy-hitters, but for others may only learn the {\em combined counts} of several values.

More concretely, each value in POPSTAR is mapped to a short bit string $\ell$, over which the server builds a prefix tree, revealing nodes in the tree as the counts of their parent node exceed the threshold $t$.
At the start of aggregation, the server only learns the first bit of $\ell$ for each report.
Then, if at least $t$ reports have the prefix $0$, the server learns the two bit prefix for all of those reports, and so on for all other prefixes that meet the threshold.
When a given prefix $\ell$ has been fully revealed, STAR-like tags for reports with that prefix are revealed in turn, allowing the Aggregation Server to recover values as in STAR.
The consequences of this are that some prefixes may not have over $t$ reports associated with them, and therefore do not reveal how many values were mapped to that prefix.
\cite{\citepopstar} refers to these prefixes and the data associated with them as \emph{end nodes}.
Prefixes with at least $t$ associated reports reveal the \emph{exact counts} of non-heavy-hitter values, as they return to a STAR-like tagging scheme.

The protocol described in this work introduces a different leakage regime to the one considered in POPSTAR. Rather than leak {\em some} information about {\em every} report, we leak {\em complete} information (including the underlying value) on {\em some} reports.
This leakage stems from the MDSS ramp.
When reports are points on polynomials of degree $\pd$, MDSS guarantees that a value can be recovered when there are at least $\threshold \gtrsim \pd + \frac{\max}{\np}$ reports for it among a pool of $N$ reports, where $\np$ is a parameter configured at setup time.
Therefore, we leak all values reported between $\pd$ and $\pd + \frac{\max}{\np}$ times: they are potentially recoverable as the number of reports is greater than the degree of the polynomials, but we cannot guarantee their recovery and so they are not considered heavy-hitters.

For our first experiment in \hme{ADD BACK}
% \cref{sec:mdss-experiments} 
we set $\np$ so that the number of values leaked in our scheme is the same in expectation as the number of exact counts leaked by POPSTAR.
For comparisons to other private aggregation systems, see \hme{ADD BACK}.
% \cref{sec:related-works}.
